\documentclass{article}

\usepackage{amsmath}
\usepackage{amssymb}
\usepackage{geometry}
\usepackage{csquotes}

\geometry{margin=2.5cm}
\title{Collection of Pseudo-Code definitions for the Tensor Omics
  Implementation}
\author{Asis Hallab, ``el Bobito''}

\begin{document}

\maketitle

\section{F42 Core Sorting Module Pseudocode}

\subsection{Module Overview}

The sorting module provides type-specific, in-place sorting routines for 1D arrays of integers, real numbers, and character sequences.
It uses Fortran \texttt{interface} blocks to overload a common sort function name across types.
All working memory is preallocated and passed explicitly.

\subsection{Type Definitions}

\begin{itemize}
    \item No new types are defined.
    \item Only standard Fortran arrays are used.
\end{itemize}

\subsection{Public Interface}

\begin{itemize}
    \item \texttt{interface sort\_array}
    \begin{itemize}
        \item \texttt{module procedure sort\_integer, sort\_real, sort\_character}
    \end{itemize}
\end{itemize}

Each sort routine works in-place and expects preallocated working memory.

\subsection{Function Signatures}

\begin{itemize}
    \item \texttt{subroutine sort\_integer(array, work)}
    \begin{itemize}
        \item \texttt{integer, dimension(n), intent(inout) :: array}
        \item \texttt{integer, dimension(n), intent(inout) :: work} \hfill (temporary work array)
    \end{itemize}

    \item \texttt{subroutine sort\_real(array, work)}
    \begin{itemize}
        \item \texttt{real, dimension(n), intent(inout) :: array}
        \item \texttt{integer, dimension(n), intent(inout) :: work}
    \end{itemize}

    \item \texttt{subroutine sort\_character(array, work)}
    \begin{itemize}
        \item \texttt{character(len=*), dimension(n), intent(inout) :: array}
        \item \texttt{integer, dimension(n), intent(inout) :: work}
    \end{itemize}
\end{itemize}

\subsection{Algorithm: General Structure (All Types)}

\begin{enumerate}
    \item Initialize permutation array:
    \[
    \text{work}(i) = i \quad \text{for} \quad i = 1, \ldots, n
    \]
    \item Sort the permutation array \texttt{work} based on the corresponding values in \texttt{array}.
    \begin{itemize}
        \item Use insertion sort (small \( n \)) or quicksort (large \( n \)).
    \end{itemize}
    \item Rearrange \texttt{array} according to the sorted permutation:
    \begin{itemize}
        \item Use a cycle leader algorithm to avoid extra memory allocation.
    \end{itemize}
\end{enumerate}

\subsection{Details for Each Type}

\begin{itemize}
    \item \textbf{Integer}: Compare \texttt{array(i)} and \texttt{array(j)} numerically.
    \item \textbf{Real}: Compare \texttt{array(i)} and \texttt{array(j)} numerically.
    \item \textbf{Character}: Compare \texttt{array(i)} and \texttt{array(j)} lexicographically.
\end{itemize}

\subsection{Memory Requirements}

\begin{itemize}
    \item For sorting an array of size \( n \):
    \begin{itemize}
        \item One integer work array of size \( n \) is needed.
    \end{itemize}
\end{itemize}

\subsection{Algorithm: Non-Recursive Quicksort for Permutation Array}

To comply with F42 and WebAssembly restrictions, Quicksort is implemented without recursion,  
using an explicit manual stack structure for subarray boundaries.

\subsubsection{Procedure: \texttt{quicksort\_nonrecursive}}

\begin{itemize}
    \item \texttt{subroutine quicksort\_nonrecursive(array, perm, n, stack\_left, stack\_right)}
    \begin{itemize}
        \item \texttt{type depends on context :: array} \hfill (integer, real, or character array)
        \item \texttt{integer, dimension(n), intent(inout) :: perm} \hfill (permutation array)
        \item \texttt{integer, intent(in) :: n} \hfill (size of array)
        \item \texttt{integer, dimension(:), intent(inout) :: stack\_left, stack\_right} \hfill (manual stack arrays)
    \end{itemize}
\end{itemize}

\subsubsection{Working Memory Requirements}

\begin{itemize}
    \item \texttt{stack\_left} and \texttt{stack\_right}: integer arrays of length at least \(\log_2(n) + 10\) (safe overestimate).
\end{itemize}

\subsubsection{Steps}

\begin{enumerate}
    \item Initialize:
    \begin{itemize}
        \item Set \texttt{stack\_ptr = 1}.
        \item Push initial full range: \texttt{stack\_left(1) = 1}, \texttt{stack\_right(1) = n}.
    \end{itemize}
    
    \item Main loop:
    \begin{itemize}
        \item While \texttt{stack\_ptr > 0}:
        \begin{enumerate}
            \item Pop:
            \begin{itemize}
                \item Set \texttt{left = stack\_left(stack\_ptr)}, \texttt{right = stack\_right(stack\_ptr)}.
                \item Decrement \texttt{stack\_ptr}.
            \end{itemize}
            \item If \texttt{left} \(\geq\) \texttt{right}, continue to next loop iteration.
            
            \item Select pivot:
            \begin{itemize}
                \item Set \texttt{pivot\_index = (left + right) / 2}.
                \item Set \texttt{pivot\_value = array(perm(pivot\_index))}.
            \end{itemize}
            
            \item Partition:
            \begin{itemize}
                \item Initialize \texttt{i = left}, \texttt{j = right}.
                \item While \texttt{i $\leq$ j}:
                \begin{itemize}
                    \item Increment \texttt{i} while \texttt{array(perm(i)) < pivot\_value}.
                    \item Decrement \texttt{j} while \texttt{array(perm(j)) > pivot\_value}.
                    \item If \texttt{i} \(\leq\) \texttt{j}:
                    \begin{itemize}
                        \item Swap \texttt{perm(i)} and \texttt{perm(j)}.
                        \item Increment \texttt{i}, decrement \texttt{j}.
                    \end{itemize}
                \end{itemize}
            \end{itemize}
            
            \item Push new ranges onto the stack:
            \begin{itemize}
                \item If \texttt{left < j}:
                \begin{itemize}
                    \item Increment \texttt{stack\_ptr}.
                    \item Set \texttt{stack\_left(stack\_ptr) = left}, \texttt{stack\_right(stack\_ptr) = j}.
                \end{itemize}
                \item If \texttt{i < right}:
                \begin{itemize}
                    \item Increment \texttt{stack\_ptr}.
                    \item Set \texttt{stack\_left(stack\_ptr) = i}, \texttt{stack\_right(stack\_ptr) = right}.
                \end{itemize}
            \end{itemize}
        \end{enumerate}
    \end{itemize}
\end{enumerate}

\subsubsection{Notes}

\begin{itemize}
    \item No recursion: stack is manually managed.
    \item Memory predictable: all working memory passed as arguments.
    \item Stateless: no hidden or module-wide variables.
    \item Safe for WebAssembly compilation and fully F42 compliant.
\end{itemize}

\section{Unit Tests for F42 Core Sorting Module}

The following unit tests verify correctness, robustness, and F42 compliance of the non-recursive sorting routines for integer, real, and character arrays.
Each test assumes that all required working memory is preallocated before calling sort procedures.

\subsection{Test 1: Sort Ascending Integer Array}

\begin{itemize}
    \item Input:
    \begin{itemize}
        \item \texttt{array} = [5, 2, 9, 1, 6]
        \item \texttt{work} = preallocated integer array of size 5
    \end{itemize}
    \item Action:
    \begin{itemize}
        \item Call \texttt{sort\_array(array, work)}.
    \end{itemize}
    \item Expected:
    \begin{itemize}
        \item \texttt{array} = [1, 2, 5, 6, 9]
    \end{itemize}
\end{itemize}

\subsection{Test 2: Sort Descending Real Array}

\begin{itemize}
    \item Input:
    \begin{itemize}
        \item \texttt{array} = [3.5, 2.2, 8.8, 1.1, 7.7]
        \item \texttt{work} = preallocated integer array of size 5
    \end{itemize}
    \item Action:
    \begin{itemize}
        \item Call \texttt{sort\_array(array, work)}.
    \end{itemize}
    \item Expected:
    \begin{itemize}
        \item \texttt{array} = [1.1, 2.2, 3.5, 7.7, 8.8]
    \end{itemize}
\end{itemize}

\subsection{Test 3: Sort Random Character Array}

\begin{itemize}
    \item Input:
    \begin{itemize}
        \item \texttt{array} = ["dog", "apple", "zebra", "cat", "bird"]
        \item \texttt{work} = preallocated integer array of size 5
    \end{itemize}
    \item Action:
    \begin{itemize}
        \item Call \texttt{sort\_array(array, work)}.
    \end{itemize}
    \item Expected:
    \begin{itemize}
        \item \texttt{array} = ["apple", "bird", "cat", "dog", "zebra"]
    \end{itemize}
\end{itemize}

\subsection{Test 4: Stability on Already Sorted Arrays}

\begin{itemize}
    \item Input:
    \begin{itemize}
        \item \texttt{array} = [1, 2, 3, 4, 5]
        \item \texttt{work} = preallocated integer array of size 5
    \end{itemize}
    \item Action:
    \begin{itemize}
        \item Call \texttt{sort\_array(array, work)}.
    \end{itemize}
    \item Expected:
    \begin{itemize}
        \item \texttt{array} remains [1, 2, 3, 4, 5]
    \end{itemize}
\end{itemize}

\subsection{Test 5: Behavior on Empty Array}

\begin{itemize}
    \item Input:
    \begin{itemize}
        \item \texttt{array} = []
        \item \texttt{work} = preallocated integer array of size 0
    \end{itemize}
    \item Action:
    \begin{itemize}
        \item Call \texttt{sort\_array(array, work)}.
    \end{itemize}
    \item Expected:
    \begin{itemize}
        \item No error or crash; operation completes trivially.
    \end{itemize}
\end{itemize}

\subsection{Test 6: Correctness of Sorting After Random Shuffles}

\begin{itemize}
    \item Input:
    \begin{itemize}
        \item Generate random permutation of integers [1,2,3,...,1000].
        \item \texttt{array} = shuffled integers
        \item \texttt{work} = preallocated integer array of size 1000
    \end{itemize}
    \item Action:
    \begin{itemize}
        \item Call \texttt{sort\_array(array, work)}.
    \end{itemize}
    \item Expected:
    \begin{itemize}
        \item \texttt{array} = [1, 2, 3, ..., 1000]
    \end{itemize}
\end{itemize}


\section{F42 Adaptive Gaussian k-Nearest Neighbor Smoothing (``Marduk'')}

\begin{displayquote}
In Assyrian mythology, \emph{Marduk} is known as the ``placer of stars''
due to his role in creating the heavens and earth, including the placement of
constellations and the stars themselves.
\end{displayquote}

\subsection{Vector-Valued Smoothing Routine \texttt{adapt\_gauss\_knn\_smooth}}

This routine performs local smoothing of a vector-valued field (e.g., expression or shift vectors) by computing a Gaussian-weighted average over the $k$ nearest neighbors of a given target point. The kernel width $\sigma$ is adapted for each target point based on the distance to its $k$-th nearest neighbor. The smoothed vector is written to the output \texttt{result}.

\subsubsection*{Function Signature}

\begin{verbatim}
subroutine adapt_gauss_knn_smooth(
    vecs,             ! real(ndim, n_points)
    inds,             ! logical(n_points)
    target_index,     ! integer
    k,                ! integer
    kd_tree,          ! prebuilt K-D tree structure
    scratch_indices,  ! integer(k)
    scratch_dists,    ! real(k)
    result            ! real(ndim)
)
\end{verbatim}

\subsubsection*{Argument Explanation}

\begin{itemize}
  \item \texttt{vecs}: A \texttt{real(ndim, n\_points)} matrix containing the vector field to be smoothed. Each column is a vector in $\mathbb{R}^{ndim}$.
  \item \texttt{inds}: A logical mask array of length \texttt{n\_points}. Only entries with \texttt{.true.} are considered valid data points for smoothing.
  \item \texttt{target\_index}: The index in \texttt{vecs} of the point to smooth. This must be one of the \texttt{.true.} entries in \texttt{inds}.
  \item \texttt{k}: The number of nearest neighbors to use.
  \item \texttt{kd\_tree}: A precomputed K-D tree for \texttt{vecs(:, inds)}, assumed to support masked neighbor lookup.
  \item \texttt{scratch\_indices}: Transient integer array of length at least \texttt{k} for neighbor indices.
  \item \texttt{scratch\_dists}: Transient real array of length \texttt{k} for distances to the neighbors.
  \item \texttt{result}: Output array of dimension \texttt{ndim} containing the smoothed vector.
\end{itemize}

\textbf{Note}: This function produces the smoothed vector only for the single \texttt{target\_index}. To smooth an entire subfield, the function must be called iteratively for all entries with \texttt{.true.} in \texttt{inds}.

\subsubsection*{Pseudocode}

\begin{verbatim}
1. target_coords := vecs(:, target_index)

2. call find_k_nearest_neighbors_masked(
       kd_tree,
       target_coords,
       inds,
       k,
       scratch_indices,
       scratch_dists
   )

3. d_max := scratch_dists(k)
   sigma := d_max / 3.0

4. result := 0.0_vector(ndim)
   total_weight := 0.0

5. do i = 1, k
     dist := scratch_dists(i)
     w := exp( - (dist**2) / (2.0 * sigma**2) )

     idx := scratch_indices(i)
     result := result + w * vecs(:, idx)
     total_weight := total_weight + w
   end do

6. if total_weight > 0.0 then
     result := result / total_weight
   else
     result := vecs(:, target_index)  ! fallback: no smoothing
   end if
\end{verbatim}

\subsubsection*{Usage Pattern for Full Subfield Smoothing}

To smooth all vectors in a subfield defined by \texttt{inds}, iterate over all \texttt{target\_index} such that \texttt{inds(target\_index) = .true.}, and call \texttt{adapt\_gauss\_knn\_smooth} for each.

This guarantees full F42 compatibility, vector-wise smoothing, and modular use with Tensor Omics vector field data structures.

\subsection{Applications Across TOX}

\begin{itemize}
  \item Normalization: Fit for gene-wise SD vs. mean.
  \item Interpolation and smoothing of expression vectors or shift vectors.
  \item Lexical smoothing across metadata axes.
  \item Clock hand and trajectory smoothing (e.g. Mjölnir-aligned).
  \item Local field smoothing in GAWD or Bobito cones.
\end{itemize}

\subsection{Test Suite for \texttt{adapt\_gauss\_knn\_smooth}}

This suite verifies the correctness of the adaptive Gaussian k-nearest neighbor smoothing algorithm across edge cases, synthetic inputs, and known outcomes.

\subsubsection*{Test 1: Single Point, Self Only}

\textbf{Goal}: Verify that smoothing a vector when only one point is available returns that point unchanged.

\begin{verbatim}
Input:
  vecs := [1.0, 2.0, 3.0]  ! ndim = 3, n_points = 1
  inds := [.true.]
  target_index := 1
  k := 1
  result := [0.0, 0.0, 0.0]

Expected Output:
  result == [1.0, 2.0, 3.0]
\end{verbatim}

---

\subsubsection*{Test 2: Two Identical Vectors, k = 2}

\textbf{Goal}: Check that smoothing over two identical vectors returns the same vector.

\begin{verbatim}
Input:
  vecs := [[1.0, 1.0],
           [2.0, 2.0],
           [3.0, 3.0]]  ! ndim = 3, n_points = 2
  inds := [.true., .true.]
  target_index := 1
  k := 2

Expected Output:
  result == [1.0, 2.0, 3.0]  ! unchanged because all neighbors identical
\end{verbatim}

---

\subsubsection*{Test 3: Linear Interpolation Between Two Vectors}

\textbf{Goal}: Validate that Marduk performs linear interpolation for symmetric neighbors.

\begin{verbatim}
Input:
  vecs := [[0.0, 2.0],
           [0.0, 2.0],
           [0.0, 2.0]]  ! ndim = 3, n_points = 2
  inds := [.true., .true.]
  target_index := 1
  k := 2

Expected Output:
  result approximately equals [1.0, 1.0, 1.0]  ! mean of two equidistant neighbors
\end{verbatim}

---

\subsubsection*{Test 4: Gaussian Weight Decay Check}

\textbf{Goal}: Ensure closer neighbor contributes more via Gaussian kernel.

\begin{verbatim}
Input:
  vecs := [[0.0, 10.0],
           [0.0, 10.0],
           [0.0, 10.0]]
  inds := [.true., .true.]
  target_index := 1
  k := 2

Expected Output:
  result closer to [0.0, 0.0, 0.0] than to [10.0, 10.0, 10.0]
\end{verbatim}

---

\subsubsection*{Test 5: Masked Points Exclusion}

\textbf{Goal}: Confirm that \texttt{.false.} entries in \texttt{inds} are never used.

\begin{verbatim}
Input:
  vecs := [[0.0, 10.0, 20.0],
           [0.0, 10.0, 20.0],
           [0.0, 10.0, 20.0]]
  inds := [.true., .false., .true.]
  target_index := 1
  k := 2

Expected Output:
  result is mean of point 1 and point 3, not including masked-out point 2
\end{verbatim}

---

\subsubsection*{Test 6: Zero Neighbors Edge Case}

\textbf{Goal}: Verify fallback behavior if no valid neighbors found.

\begin{verbatim}
Input:
  vecs := [[1.0], [2.0], [3.0]]
  inds := [.false.]  ! target not in mask
  target_index := 1
  k := 1

Expected Output:
  result == vecs(:, 1)  ! fallback to unsmoothed value
\end{verbatim}

---

\subsubsection*{Test 7: Dimensional Scaling Check}

\textbf{Goal}: Confirm algorithm generalizes to arbitrary ndim.

\begin{verbatim}
Input:
  vecs := real(10, 100)  ! ndim = 10, n_points = 100
  Populate with smooth gradient vectors
  inds := all .true.
  target_index := random index
  k := 10

Expected Output:
  result approximately equals smooth interpolation within local 10D neighborhood
\end{verbatim}


\subsection{The \texttt{which} Function for Logical Arrays (F42-Compliant)}

The \texttt{which} function extracts the indices of all entries in a logical array that evaluate to \texttt{.TRUE.}. This routine is the direct analogue of the R function \texttt{which()}, designed to work with logical arrays resulting from any Fortran logical expression (e.g., \texttt{x > 3.0 .AND. MOD(x,2) == 0}). It follows the F42 principles: stateless, allocation-free, and usable in high-performance settings.

\subsubsection{Interface Specification}

\begin{verbatim}
subroutine which(
    mask(n),            ! Input logical array
    n,                  ! Size of the input array
    idx_out(m_max),     ! Preallocated output array of indices
    m_max,              ! Capacity of idx_out
    m_out               ! Number of matches returned
)
\end{verbatim}

\begin{itemize}
  \item \texttt{mask(n)}: Logical array derived from a Fortran logical expression
  \item \texttt{idx\_out(m\_max)}: Output array holding 1-based indices of matches
  \item \texttt{m\_out}: Actual number of matches written to \texttt{idx\_out}
\end{itemize}

\subsubsection{Pseudocode Implementation}

\begin{verbatim}
m_out = 0
do i = 1, n
    if (mask(i)) then
        if (m_out < m_max) then
            m_out = m_out + 1
            idx_out(m_out) = i
        endif
    endif
enddo
\end{verbatim}

This implementation ensures no allocation and full compatibility with high-performance computing contexts (e.g., OpenMP, SIMD, or GPU-friendly flattening).

\subsubsection{Test Case 1: Simple Threshold Filter}

\begin{itemize}
  \item \textbf{Input:} \texttt{x = [1.0, 4.0, 6.0, 2.0, 8.0]}
  \item \textbf{Logical Mask:} \texttt{mask = (x > 3.0)}
  \item \textbf{Expected Output:} \texttt{idx\_out = [2, 3, 5]}, \texttt{m\_out = 3}
\end{itemize}

\subsubsection{Test Case 2: Composite Logical Condition}

\begin{itemize}
  \item \textbf{Input:} \texttt{x = [1.0, 4.0, 6.0, 2.0, 8.0]}
  \item \textbf{Mask:} \texttt{mask = (x > 3.0 .AND. MOD(INT(x),2) == 0)}
  \item \textbf{Expected Output:} \texttt{idx\_out = [2, 5]}, \texttt{m\_out = 2}
\end{itemize}

\subsubsection{Test Case 3: No Matches}

\begin{itemize}
  \item \textbf{Input:} \texttt{x = [1.0, 2.0, 3.0]}
  \item \textbf{Mask:} \texttt{mask = (x > 10.0)}
  \item \textbf{Expected Output:} \texttt{idx\_out = []}, \texttt{m\_out = 0}
\end{itemize}

\subsubsection{Performance and Compliance Notes}

\begin{itemize}
  \item This function is safe for use in any high-performance loop or parallel region.
  \item Overflow behavior (when \texttt{m\_out} exceeds \texttt{m\_max}) must be handled externally.
  \item All array bounds and allocations are managed by the caller.
\end{itemize}

\section{F42-Compliant Index-Based K-D Trees and BSTs}

\subsection{Overview}

This section defines efficient, allocation-free, recursion-free, and Fortran-compatible implementations of:
\begin{itemize}
    \item A 1D Binary Search Tree (BST) represented as a flat index array
    \item A Cartesian K-D Tree using descending variance dimension order
    \item A Spherical K-D Tree (used for directional queries based on cosine similarity)
    \item Core access and query routines
    \item Unit tests for each structure
\end{itemize}

All memory structures are passed in. The trees are represented purely as sorted index arrays. No structs, recursion, or dynamic memory is used.

\subsection{Binary Search Tree (BST) Construction via Sorting}

\textbf{Input:}
\begin{itemize}
    \item Real array \texttt{x(n)} --- the data to be indexed
    \item Integer array \texttt{ix(n)} --- index array to be filled
\end{itemize}

\textbf{Pseudocode:}
\begin{verbatim}
subroutine build_bst_index(x, n, ix)
  real, intent(in) :: x(n)
  integer, intent(out) :: ix(n)

  call quicksort_indices(x, ix, n)
end subroutine
\end{verbatim}

\subsection{Access Functions for BST (Flat Index Array)}

\paragraph{Get Value at Sorted Position}
Returns the value of the original data array at the position given by the sorted index.

\begin{verbatim}
function get_sorted_value(x, ix, i) result(val)
real, intent(in) :: x(:)
integer, intent(in) :: ix(:)
integer, intent(in) :: i
real :: val

val = x(ix(i))
end function
\end{verbatim}

\paragraph{Get Sorted Index for Range Query}
Example: Find values in \texttt{x} between a given \texttt{lo} and \texttt{hi}. This function returns a list of positions \texttt{i} such that \texttt{lo <= x(ix(i)) <= hi}.

\begin{verbatim}
subroutine bst_range_query(x, ix, n, lo, hi, out_ix, out_n)
real, intent(in) :: x(:)
integer, intent(in) :: ix(:)
integer, intent(in) :: n
real, intent(in) :: lo, hi
integer, intent(out) :: out_ix(n)
integer, intent(out) :: out_n

integer :: i
out_n = 0

do i = 1, n
if (x(ix(i)) >= lo .and. x(ix(i)) <= hi) then
out_n = out_n + 1
out_ix(out_n) = ix(i)
else if (x(ix(i)) > hi) then
exit ! Since array is sorted, we can stop early
end if
end do
end subroutine
\end{verbatim}

\paragraph{Use Case Example}
Sort an array and extract all values between 2.0 and 3.5.

\begin{verbatim}
real :: x(1000)
integer :: ix(1000), res_ix(1000), res_n

call random_fill(x)
call build_bst_index(x, 1000, ix)
call bst_range_query(x, ix, 1000, 2.0, 3.5, res_ix, res_n)
! Now x(res_ix(1:res_n)) holds the sorted matching values
\end{verbatim}

\subsection{Cartesian K-D Tree Construction}

\textbf{Input:}
\begin{itemize}
    \item Real matrix \texttt{X(d, n)} --- the data points
    \item Integer array \texttt{kd\_ix(n)} --- output index array representing K-D Tree
    \item Integer array \texttt{dim\_order(d)} --- precomputed dimension order by descending variance
    \item Integer work arrays for temporary use
\end{itemize}

\textbf{Pseudocode:}
\begin{verbatim}
subroutine build_kd_index(X, d, n, kd_ix, dim_order, work)
  ! Loop-based K-D Tree construction by sorting and partitioning
  ! No recursion. Uses stack to manage index intervals and split dimension.

  integer :: stack_top, stack(1:2, max_depth)
  integer :: l, r, dim, mid, depth, i

  stack_top = 1
  stack(:, 1) = [1, n]

  do while (stack_top > 0)
    l = stack(1, stack_top)
    r = stack(2, stack_top)
    stack_top = stack_top - 1
    if (r <= l) cycle

    depth = some_depth_heuristic(l, r)  ! Or simply count levels
    dim = dim_order(mod(depth, d) + 1)

    call partial_sort_by_dimension(X, d, kd_ix, l, r, dim)
  end do
end subroutine
\end{verbatim}

\subsection{Spherical K-D Tree for Directional Queries}

\textbf{Input:}
\begin{itemize}
    \item Normalized vector matrix \texttt{V(d, n)} --- each column is unit length
    \item Integer array \texttt{sphere\_ix(n)} --- output index array for spherical K-D Tree
\end{itemize}

\textbf{Pseudocode:}
\begin{verbatim}
subroutine build_spherical_kd(V, d, n, sphere_ix, dim_order)
  call build_kd_index(V, d, n, sphere_ix, dim_order)
end subroutine
\end{verbatim}

\subsection{Access Routines}

\textbf{Get value at sorted position:}
\begin{verbatim}
function get_value_sorted(x, ix, i)
  real :: get_value_sorted
  real, intent(in) :: x(:)
  integer, intent(in) :: ix(:)
  get_value_sorted = x(ix(i))
end function
\end{verbatim}

\textbf{Get point from KD index:}
\begin{verbatim}
subroutine get_kd_point(X, kd_ix, i, point)
  real, intent(in) :: X(:, :)
  integer, intent(in) :: kd_ix(:)
  real, intent(out) :: point(:)
  point = X(:, kd_ix(i))
end subroutine
\end{verbatim}

\subsection{Unit Tests}

\textbf{BST Test:}
\begin{verbatim}
call random_array(x, n)
call build_bst_index(x, n, ix)
assert_monotonic(x(ix))
\end{verbatim}

\textbf{K-D Tree Test:}
\begin{verbatim}
call random_matrix(X, d, n)
call compute_dim_order(X, d, n, dim_order)
call build_kd_index(X, d, n, kd_ix, dim_order)
assert_valid_kd_structure(kd_ix)
\end{verbatim}

\end{document}
